\documentclass[11pt,a4paper]{article}
\usepackage[utf8]{inputenc}
\usepackage[T1]{fontenc}
\usepackage{geometry}
\geometry{margin=2.5cm}
\usepackage{graphicx}
\usepackage{booktabs}
\usepackage{array}
\usepackage{xcolor}
\usepackage{listings}
\usepackage{hyperref}
\usepackage{float}
\usepackage{caption}
\usepackage{multirow}
\usepackage{enumitem}

\lstdefinestyle{Bits}{% bit-field diagrams: no line numbers, small enough to fit
  language={},basicstyle=\ttfamily\scriptsize,numbers=none,breaklines=false,
  frame=none,showstringspaces=false,columns=fullflexible,keepspaces=true,
  aboveskip=6pt,belowskip=10pt,xleftmargin=2pt
}
\lstdefinestyle{Cstyle}{
  language=C,
  basicstyle=\ttfamily\footnotesize,
  numbers=left,
  numberstyle=\tiny,
  frame=single,
  breaklines=true,
  tabsize=2,
  commentstyle=\color{gray},
  keywordstyle=\color{blue},
  stringstyle=\color{red},
}
\lstdefinestyle{asmstyle}{
  language={[x86masm]Assembler},
  basicstyle=\ttfamily\footnotesize,
  numbers=left,
  numberstyle=\tiny,
  frame=single,
  breaklines=true,
  tabsize=2,
  commentstyle=\color{gray},
}

\hypersetup{
  colorlinks=true,
  pdftitle={BradISA V2 Specification},
  pdfauthor={Brad Devices},
}

\usepackage{bradbrand}

\begin{document}
\bradtitlepage{BradISA V2}{%
  \textbf{BradISA V2}}{%
  Instruction Set Architecture Specification}{%
  Revision 1.0 --- Kinetic Gen1 --- July 2026}
\newpage

\tableofcontents
\newpage

\section{Introduction}

BradISA V2 (Kinetic Gen2) is a 32-bit fixed-length RISC instruction set
architecture designed for the Brad Devices BradCore processor family. It
provides 32 general-purpose registers, 64 opcodes encoded in a 6-bit primary
opcode field, and seven instruction formats, each of which is exactly 32 bits
wide.

V2 supersedes V1 in full. It is a clean break: every instruction encoding
changes, and none of V1's encodings are preserved. That is a deliberate
consequence of the field widths --- growing a register field from four bits to
five shifts every field above it, so no compatibility-preserving arrangement
exists once the register file grows. The cost is nil today, because no V1
binary has ever executed on silicon; the same break after a toolchain ships
would be a migration.

\subsection{Why V2 exists}

V1 cannot be compiled. Its RI format carries a 16-bit sign-extended immediate
and no instruction loads a wide constant, so the addressable range from any base
register is $\pm$32\,KiB. There is therefore no data segment beyond 64\,KiB, no
heap, no array longer than 16,384 elements and no linked executable --- which
means no C library. V1 has no indirect jump, so no function pointers, no
virtual calls and no callbacks.

The three missing instruction groups are a symptom rather than the disease.
V1's sixteen opcodes are fully allocated, \texttt{0x0} \texttt{ADD} through
\texttt{0xF} \texttt{RET}, with none free. The property the V1 design document held
up as its central argument --- a sixteen-case scalar decoder with no decode
hierarchy --- is precisely the property that prevents it growing. An ISA cannot
have both a small exhaustive opcode space and the ability to address memory.

\subsection{Why 64 opcodes rather than 32}

An earlier proposal in this review recommended a 5-bit field, 32 opcodes. The
arithmetic does not support it. Clang's integer subset alone requires:

\begin{table}[htbp]
\centering
\caption{Integer instructions a compiler emits, against opcode-space size}
\begin{tabular}{lp{8.6cm}r}
\toprule
Group & Instructions & Count \\
\midrule
ALU, register-register   & ADD SUB MUL MULHU DIV DIVU REM REMU AND OR XOR SHL SHR SAR & 12 \\
ALU, register-immediate  & ADDI ANDI ORI XORI SLLI SRLI SRAI SLTI SLTIU & 9 \\
Load / store             & LDW LHU LHBU STW SH SB & 6 \\
Upper immediate          & LUI AUIPC & 2 \\
Control transfer         & NOP JMP JAL CALL JALR RET FENCE & 6 \\
Branches                 & BEQ BNE BLT BGE BLTU BGEU BZ BNZ & 8 \\
Atomics                  & LR.W SC.W AMOSWAP.W AMOADD.W & 4 \\
\midrule
Total                    &  & 50 \\
\bottomrule
\end{tabular}
\end{table}

A 32-slot space is eleven slots short of the fifty instructions above. Closing
that gap at 5 bits requires either folding the register-immediate ALU variants
behind a \texttt{funct} field --- reintroducing the sub-hierarchy the V1 design
explicitly rejected --- or deferring atomics to a third encoding break. V2 takes
neither cost and spends one more bit.

\begin{table}[htbp]
\centering
\caption{Opcode-space arithmetic}
\begin{tabular}{lll}
\toprule
Field width & Slots & Verdict \\
\midrule
4 bits & 16 & V1. Short by 34. \\
5 bits & 32 & Short by 18 against 50. \\
6 bits & 64 & Fits, with 14 slots reserved and accounted for. \\
\bottomrule
\end{tabular}
\end{table}

\subsection{Why six bits is not new to this project}

The vector core already implements this encoding.

\begin{lstlisting}[style=Cstyle]
  src/rtl/bradvector_core.v:27   localparam REG_ADDR_W = 5
  src/rtl/bradvector_core.v:28   localparam NUM_REGS     = 32
  src/rtl/bradvector_core.v:41   wire [4:0] slot0_opcode
\end{lstlisting}

The scalar ISA was the outlier that never grew. V2 ports the vector core's
proven field widths to the scalar side, so the argument for the change is
consistency between two cores already in the same repository rather than a new
mechanism.

\subsection{Data Model}
\begin{itemize}[nosep]
  \item \textbf{Word size:} 32 bits (4 bytes)
  \item \textbf{Address space:} 32-bit byte-addressed
  \item \textbf{Page size:} 4096 bytes (4\,KiB)
  \item \textbf{Endianness:} Little-endian
  \item \textbf{Instruction length:} 32 bits, fixed
\end{itemize}

\subsection{Data Model}
\begin{itemize}[nosep]
  \item \textbf{Word size:} 32 bits (4 bytes)
  \item \textbf{Address space:} 32-bit byte-addressed
  \item \textbf{Page size:} 4096 bytes (4\,KiB)
  \item \textbf{Endianness:} Little-endian
  \item \textbf{Instruction length:} 32 bits, fixed
\end{itemize}

\subsection{Pipeline Models}
BradISA defines two pipeline models:

\begin{table}[htbp]
\centering
\caption{Pipeline model comparison}
\begin{tabular}{lcc}
\toprule
Property               & Phoenix            & Falcon           \\
\midrule
Pipeline depth         & 10 stages          & 5 stages         \\
Issue width            & Octa-issue         & Single-issue     \\
Execution order        & Out-of-order       & In-order         \\
ROB size              & 128 entries         & ---              \\
Typical IPC           & 2.2                & 0.9              \\
Peak IPC              & 5.0                & 1.0              \\
Target DMIPS/MHz      & 5.5                & 2.2              \\
Target frequency      & 3.0--5.2\,GHz      & 2.4--3.6\,GHz    \\
\bottomrule
\end{tabular}
\end{table}

\newpage

\section{Register File}

BradISA V2 provides 32~32-bit general-purpose registers, numbered {\tt r0}
through {\tt r31}. Register {\tt r0} is hardwired to zero; writes to {\tt r0}
are silently discarded.

\subsection{Why V2 grew the register file}

V1 provided sixteen registers. Widening the register fields from four bits to
five is the same one-bit change as widening the opcode, and it brings the
scalar core into line with the vector core, which already ships
\texttt{REG\_ADDR\_W = 5} and \texttt{NUM\_REGS = 32}.

The motivation is not the field width. It is the ABI. V1 allocated its sixteen
registers as four arguments, three temporaries, five callee-saved, a stack
pointer, a link register and a read-only program counter, which left a register
allocator roughly eight usable registers and the calling convention four
parameters. Every call past the fourth spilled. A compiler can be written for
such an ABI and it will work, but it will spill continuously, and no backend
quality survives it.

V2 settles this in \S\ref{sec:abi}: eight argument registers, six temporaries,
eleven callee-saved, and no architectural PC register.

\subsection{ABI Convention}
\label{sec:abi}

\textbf{Decision.} V2 passes arguments in eight registers, holds six
temporaries and eleven callee-saved, and \textbf{retires the architectural PC
register}. The shape is x86-32's argument count on a RISC-V register budget.
The PC pseudo-register occupied one of thirty-two slots while being
unencodable as {\tt rd}, {\tt rs1} or {\tt rs2}; removing it costs nothing,
because {\tt AUIPC} already supplies PC-relative addressing and an explicit
``read the PC'' operation can take a reserved opcode when one is wanted.

\begin{table}[htbp]
\centering
\caption{Application Binary Interface register map (32 registers)}
\begin{tabular}{cllp{5.4cm}}
\toprule
Register & ABI Name & Count & Purpose \\
\midrule
{\tt r0}       & {\tt zero} & 1  & Hardwired zero; writes discarded \\
{\tt r1}--{\tt r8}   & {\tt a0}--{\tt a7} & 8  & Arguments~\#0--\#7; {\tt a0} also returns \\
{\tt r9}--{\tt r14}  & {\tt t0}--{\tt t5} & 6  & Caller-saved temporaries \\
{\tt r15}      & {\tt lr}  & 1  & Link register; return address on {\tt CALL} \\
{\tt r16}--{\tt r26} & {\tt s0}--{\tt s10} & 11 & Callee-saved \\
{\tt r27}      & {\tt sp}  & 1  & Stack pointer \\
{\tt r28}      & {\tt gp}  & 1  & Global pointer \\
{\tt r29}      & {\tt tp}  & 1  & Thread pointer \\
{\tt r30}--{\tt r31} & ---   & 2  & Reserved \\
\bottomrule
\end{tabular}
\end{table}

Twenty-nine of thirty-two registers are allocatable: one is hardwired zero and
two are reserved for growth. The arithmetic is checked rather than asserted --
8 arguments + 6 temporaries + 11 callee-saved + \texttt{lr} + \texttt{sp} +
\texttt{gp} + \texttt{tp} = 29.

\subsection{Why this shape}

V1 allocated sixteen registers as four arguments, three temporaries, five
callee-saved, a stack pointer, a link register and a read-only program
counter. That left a register allocator roughly eight usable registers and the
calling convention four parameters, so every call past the fourth spilled to
the stack. A compiler can be written for that ABI and it will work --- it will
also spill continuously, and no backend quality survives it.

The alternative candidate was RISC-V's seven arguments. Eight was chosen
because the register file can afford it, because eight is the better-trodden
count for a 32-bit integer ABI, and because every argument beyond the register
set costs a store and a reload at both ends of every call.

Sixteen callee-saved and reserved slots remain unallocated on purpose. They are
where a vector calling convention, a software-pipeline temporaries group and
any future banked or rotate arrangement would go, and they are named as
reserved rather than left unassigned.

\subsection{Decision record}

The encoding was fixed before this convention was chosen, deliberately. Growing
a register field from four bits to five shifts every field above it, so the
encoding had to be settled first, and the convention is the part that could be
decided afterwards on evidence. It cost nothing to leave open for exactly that
length of time.

Two open questions became this one. How many arguments --- seven or eight ---
and whether to keep a read-only PC register. Both are answered above.

\subsection{Calling Convention}
\begin{itemize}[nosep]
  \item Arguments 0--7 are passed in {\tt r1}--{\tt r8}. Additional arguments are
        passed on the stack, right to left.
  \item Return value is placed in {\tt a0} ({\tt r1}).
  \item Callee-saved registers {\tt s0}--{\tt s10} must be preserved across
        function calls. {\tt t0}--{\tt t5} and {\tt a0}--{\tt a7} may be
        clobbered.
  \item The link register {\tt lr} holds the return address on {\tt CALL} and
        {\tt RET}.
  \item The stack grows downward; {\tt sp} must point to the last used stack
        slot, and the stack is 16-byte aligned on entry.
  \item There is no red zone below {\tt sp}, and no leaf-function fast path that
        omits the frame setup; a signal or exception may overwrite the
        caller-saved area.
\end{itemize}

\newpage

\section{Instruction Formats}

All BradISA V2 instructions are 32~bits wide. Seven formats are defined, and
every one of them sums to exactly 32 bits. That is not a coincidence to be
checked by hand later; it is the property the V1 specification failed to hold,
see \S\ref{sec:v1-encoding-defects}.

\begin{table}[htbp]
\centering
\caption{Instruction format field layout}
\begin{tabular}{clllllll}
\toprule
Format  & [31:26] & [25:21] & [20:16] & [15:11] & [10:6] & [5:0] \\
\midrule
{\tt RRR} & op & rd & rs1 & rs2 & --- & --- \\
{\tt RI}  & op & rd & rs1 & \multicolumn{4}{c}{imm[15:0]} \\
{\tt SI}  & op & rs2 & rs1 & \multicolumn{4}{c}{imm[15:0]} \\
{\tt BR}  & op & rs1 & rs2 & \multicolumn{4}{c}{offset[15:0]} \\
{\tt IJ}  & op & rd & rs1 & \multicolumn{4}{c}{imm[15:0]} \\
{\tt JMP} & op & rd & \multicolumn{5}{c}{offset[20:0]} \\
{\tt U}   & op & rd & \multicolumn{5}{c}{imm[20:0]} \\
\bottomrule
\end{tabular}
\end{table}

\subsection{RRR --- Register-to-Register ALU}

\begin{lstlisting}[style=Bits]
 +--------+-------+-------+-------+--------+
 | op     | rd    | rs1   | rs2   | 0      |
 +--------+-------+-------+-------+--------+   RRR
\end{lstlisting}

\[
\texttt{instruction[31:0]} = \{\, \texttt{op[5:0]}, \texttt{rd[4:0]},
  \texttt{rs1[4:0]}, \texttt{rs2[4:0]}, 11\textrm{\textquoteright}b0 \,\}
\]

Bits [10:0] are reserved and must be zero. A non-zero value in a reserved field
is an illegal instruction, not a synonym for zero, so that encodings do not
acquire two spellings.

\subsection{RI --- Register-Immediate}

Used by ADDI, ANDI, ORI, XORI, the shift-immediates, the set-less-thans and
every load.

\begin{lstlisting}[style=Bits]
 +--------+-------+-------+--------------+
 | op     | rd    | rs1   | imm          |
 +--------+-------+-------+--------------+   RI
\end{lstlisting}

\[
\texttt{instruction[31:0]} = \{\, \texttt{op[5:0]}, \texttt{rd[4:0]},
  \texttt{rs1[4:0]}, \texttt{imm[15:0]} \,\}
\]

The 16-bit immediate is sign-extended to 32 bits before use, except for
\texttt{ANDI}, \texttt{ORI} and \texttt{XORI}, which zero-extend it, because
those are bitwise operations and sign extension would be a different operation.

\subsection{SI --- Store-Immediate}

Identical bit layout to RI. The register field at [25:21] is the \emph{source}
of the stored value rather than a destination.

\begin{lstlisting}[style=Bits]
 +--------+-------+-------+--------------+
 | op     | rs2   | rs1   | imm          |
 +--------+-------+-------+--------------+   SI
\end{lstlisting}

\subsection{BR --- Conditional Branch}

\begin{lstlisting}[style=Bits]
 +--------+-------+-------+--------------+
 | op     | rs1   | rs2   | offset       |
 +--------+-------+-------+--------------+   BR
\end{lstlisting}

\[
\texttt{PC} = \texttt{insn\_pc} + 4 + \big(\texttt{sign\_extend(offset[15:0])} \times 4\big)
\]

Branch offsets are therefore $\pm$128\,KiB from the instruction. \texttt{BZ} and
\texttt{BNZ} use this format with \texttt{rs2} = 0, which is why they need no
separate encoding.

\subsection{IJ --- Indirect Jump}

\begin{lstlisting}[style=Bits]
 +--------+-------+-------+--------------+
 | op     | rd    | rs1   | imm          |
 +--------+-------+-------+--------------+   IJ
\end{lstlisting}

\texttt{JALR}: the target is computed and the link written before the transfer,
so that \texttt{jalr ra, ra, 0} is a correct tail call.

\subsection{JMP --- Unconditional Jump and Call}

\begin{lstlisting}[style=Bits]
 +--------+-------+----------------+
 | op     | rd    | offset         |
 +--------+-------+----------------+   JMP
\end{lstlisting}

\[
\texttt{PC} = \texttt{insn\_pc} + 4 + \big(\texttt{sign\_extend(offset[20:0])} \times 4\big)
\]

Unconditional jump range is $\pm$4\,MiB, unchanged from V1. The V1 encoding
spent eight bits on padding; V2 spends five of them on a real \texttt{rd}, which
is what \texttt{JAL} needs to return an address.

\subsection{U --- Upper Immediate}

\begin{lstlisting}[style=Bits]
 +--------+-------+----------------+
 | op     | rd    | imm            |
 +--------+-------+----------------+   U
\end{lstlisting}

\[
\texttt{LUI}: \quad \texttt{rd} = \texttt{sign\_extend(imm[20:0])} \ll 11
\]
\[
\texttt{AUIPC}: \quad \texttt{rd} = \texttt{PC} + \big(\texttt{sign\_extend(imm[20:0])} \ll 11\big)
\]

21 immediate bits shifted left by 11 fills all 32 bits of the destination, so
\texttt{LUI} followed by \texttt{ADDI} constructs any 32-bit constant. This is
the instruction whose absence kept V1 from addressing memory.

\subsection{RET --- Return}

Fixed encoding, with every non-opcode bit zero:

\[
\texttt{RET} = 32\textrm{\textquoteright}h\texttt{B8000000}
\]

\texttt{RET} jumps to the address in \texttt{r15 (lr)} and does not modify
\texttt{lr}.

\newpage

\subsection{Encoding diagrams}

\begin{lstlisting}[style=Bits]
 +--------+-------+-------+-------+--------+
 | op     | rd    | rs1   | rs2   | 0      |
 +--------+-------+-------+-------+--------+   RRR

 +--------+-------+-------+--------------+
 | op     | rd    | rs1   | imm          |
 +--------+-------+-------+--------------+   RI

 +--------+-------+-------+--------------+
 | op     | rs2   | rs1   | imm          |
 +--------+-------+-------+--------------+   SI

 +--------+-------+-------+--------------+
 | op     | rs1   | rs2   | offset       |
 +--------+-------+-------+--------------+   BR
\end{lstlisting}

\begin{lstlisting}[style=Bits]
 +--------+-------+-------+--------------+
 | op     | rd    | rs1   | imm          |
 +--------+-------+-------+--------------+   IJ

 +--------+-------+----------------+
 | op     | rd    | offset         |
 +--------+-------+----------------+   JMP

 +--------+-------+----------------+
 | op     | rd    | imm            |
 +--------+-------+----------------+   U

 +--------------------------------------------------------------+
|                     0xB8000000                             |   RET
 +--------------------------------------------------------------+
\end{lstlisting}

\subsection{Two encoding defects in V1, corrected here}
\label{sec:v1-encoding-defects}

Recorded because both were load-bearing for anything transcribed from this
specification, and because a TableGen target is written directly against this
table.

\textbf{The V1 RI format did not describe a 32-bit word.} The field array
summed to 28 bits ($4 + 4 + 4 + 16$), four short, and the accompanying diagram
placed \texttt{imm[15:0]} at bits 31:16, overlapping \texttt{op}, \texttt{rd}
and \texttt{rs1}. Two representations of one format, disagreeing, neither of
them a 32-bit encoding.

\textbf{V1 opcode \texttt{0xF} carried two meanings.} The scalar table assigned
\texttt{0xF = RET}; the vector extension was specified as using opcode
\texttt{0xF} for every one of its instructions, disambiguated by a class field
in bits [27:25] with \texttt{0xF0000000} pinned as the sole \texttt{RET}
encoding. The scalar and vector decoders therefore disagreed about the same
four bits, and the entire vector extension depended on \texttt{RET} having
exactly one encoding.

V2 gives the vector unit its own opcode range (\S\ref{sec:vector-escape}) and
every format sums to 32 bits.

\newpage

\section{Instruction Set}

\begin{table}[htbp]
\centering
\caption{Complete BradISA V2 integer instruction set}
\begin{tabular}{cllp{7.4cm}}
\toprule
Opcode & Mnemonic & Format & Operation \\
\midrule
0x0 & NOP    & --- & no operation \\
0x1 & ADD    & RRR & rd = rs1 + rs2 \\
0x2 & SUB    & RRR & rd = rs1 - rs2 \\
0x3 & MUL    & RRR & rd = rs1 $\times$ rs2 (low 32 bits) \\
0x4 & MULHU  & RRR & rd = (rs1 $\times$ rs2)[63:32] unsigned \\
0x5 & DIV    & RRR & rd = rs1 / rs2 signed, div-by-zero raises an exception \\
0x6 & DIVU   & RRR & rd = rs1 / rs2 unsigned \\
0x7 & REM    & RRR & rd = rs1 \% rs2 signed \\
0x8 & REMU   & RRR & rd = rs1 \% rs2 unsigned \\
0x9 & AND    & RRR & rd = rs1 \& rs2 \\
0xA & OR     & RRR & rd = rs1 | rs2 \\
0xB & XOR    & RRR & rd = rs1 \^{} rs2 \\
0xC & SHL    & RRR & rd = rs1 $<<$ rs2[4:0] \\
0xD & SHR    & RRR & rd = rs1 $>>$ rs2[4:0] (logical) \\
0xE & SAR    & RRR & rd = rs1 $\gg_s$ rs2[4:0] (arithmetic) \\
0xF & ---    & --- & reserved \\
\midrule
0x10 & ADDI   & RI  & rd = rs1 + sext(imm) \\
0x11 & ANDI   & RI  & rd = rs1 \& zext(imm) \\
0x12 & ORI    & RI  & rd = rs1 | zext(imm) \\
0x13 & XORI   & RI  & rd = rs1 \^{} zext(imm) \\
0x14 & SLLI   & RI  & rd = rs1 $<<$ imm[4:0] \\
0x15 & SRLI   & RI  & rd = rs1 $>>$ imm[4:0] (logical) \\
0x16 & SRAI   & RI  & rd = rs1 $\gg_s$ imm[4:0] (arithmetic) \\
0x17 & SLTI   & RI  & rd = (rs1 $<$ sext(imm)) ? 1 : 0 \\
0x18 & SLTIU  & RI  & rd = (rs1 $<$ zext(imm)) ? 1 : 0 \\
0x19--0x1B & --- & --- & reserved \\
0x1C--0x1F & --- & --- & reserved \\
\midrule
0x20 & LDW    & RI  & rd = mem[rs1 + sext(imm)] \\
0x21 & LHU    & RI  & rd = zeroextend(mem16[rs1 + sext(imm)]) \\
0x22 & LHBU   & RI  & rd = zeroextend(mem8[rs1 + sext(imm)]) \\
0x23 & STW    & SI  & mem[rs1 + sext(imm)] = rs2 \\
0x24 & SH     & SI  & mem16[rs1 + sext(imm)] = rs2[15:0] \\
0x25 & SB     & SI  & mem8[rs1 + sext(imm)] = rs2[7:0] \\
0x26--0x27 & --- & --- & reserved \\
\midrule
0x28 & LUI    & U   & rd = sext(imm) $<<$ 11 \\
0x29 & AUIPC  & U   & rd = PC + (sext(imm) $<<$ 11) \\
\midrule
0x2A & JMP    & JMP & PC = insn\_pc + 4 + sext(offset $\times$ 4) \\
0x2B & JAL    & JMP & rd = insn\_pc + 4; PC = target \\
0x2C & CALL   & JMP & lr = insn\_pc + 4; PC = target \\
0x2D & JALR   & IJ  & t = insn\_pc + 4; PC = (rs1 + sext(imm)) \& $\sim$1; rd = t \\
0x2E & RET    & --- & PC = lr \\
0x2F & FENCE  & --- & ordering barrier; no operand \\
\midrule
0x30 & BEQ    & BR  & if (rs1 == rs2) branch \\
0x31 & BNE    & BR  & if (rs1 != rs2) branch \\
0x32 & BLT    & BR  & if (rs1 $<$ rs2, signed) branch \\
0x33 & BGE    & BR  & if (rs1 $\geq$ rs2, signed) branch \\
0x34 & BLTU   & BR  & if (rs1 $<$ rs2, unsigned) branch \\
0x35 & BGEU   & BR  & if (rs1 $\geq$ rs2, unsigned) branch \\
0x36 & BZ     & BR  & if (rs1 == 0) branch \quad (rs2 = 0) \\
0x37 & BNZ    & BR  & if (rs1 != 0) branch \quad (rs2 = 0) \\
\midrule
0x38 & LR.W    & AMO & rd = mem[rs1]; record reservation on rs1 \\
0x39 & SC.W    & AMO & mem[rs1] = rs2; rd = 0 on success, non-zero on failure \\
0x3A & AMOSWAP.W & AMO & t = mem[rs1]; mem[rs1] = rs2; rd = t \\
0x3B & AMOADD.W  & AMO & t = mem[rs1]; mem[rs1] = t + rs2; rd = t \\
\midrule
0x3C--0x3F & --- & --- & reserved, extension space (see \S\ref{sec:vector-escape}) \\
\bottomrule
\end{tabular}
\end{table}

\subsection{Opcode allocation and what it buys}

The 64 slots are allocated by group, and every unassigned slot is stated
rather than left blank. This is the point of the exercise: an unassigned slot
in a flat decoder is not a gap, it is an instruction word that executes as
something nobody chose.

\begin{table}[htbp]
\centering
\caption{Opcode group map}
\begin{tabular}{clcl}
\toprule
Range & Group & Slots & Spare \\
\midrule
0x00--0x0F & ALU, register-register   & 15 & 1 \\
0x10--0x1F & ALU, register-immediate & 9  & 7 \\
0x20--0x27 & Load / store             & 6  & 2 \\
0x28--0x29 & Upper immediate          & 2  & 0 \\
0x2A--0x2F & Control transfer         & 6  & 0 \\
0x30--0x37 & Branches                 & 8  & 0 \\
0x38--0x3B & Atomics                  & 4  & 0 \\
0x3C--0x3F & Extension escape         & 0  & 4 \\
\midrule
\textbf{Total} & & \textbf{50} & \textbf{14} \\
\bottomrule
\end{tabular}
\end{table}

\subsection{ALU Operations}

Register-register ALU instructions use the RRR format. Register-immediate ALU
instructions use RI. Shift amounts use only the bottom five bits of the source
(0--31), in both forms.

Division and remainder are newly defined in V2 and raise an exception on a zero
divisor rather than returning a defined value, because the alternative is a
silent wrong answer.

\subsection{Memory Operations}

{\tt LDW}, {\tt LHU} and {\tt LHBU} load into {\tt rd} from
{\tt rs1 + sext(imm)}. {\tt STW}, {\tt SH} and {\tt SB} store the low 32, 16 or 8
bits of {\tt rs2} to the same effective address, and are encoded with the SI
format, in which the register field at bits [25:21] is the \emph{source}, not a
destination.

V1 used the RI shape for stores and documented this as ``STW ignores {\tt rd}
(uses {\tt rs2} as store data)''. That is the same bit layout read with two
different meanings depending on the opcode, which is exactly the sort of thing
that costs a backend an afternoon. V2 names the two cases separately.

Word accesses must be 4-byte aligned and halfword accesses 2-byte aligned. An
unaligned access raises an exception; it is not silently split.

\subsection{Addressing: how V2 reaches memory}

This is the change that makes a compiler possible.

In V1 the only addressable range from any base register was
$\pm$32\,KiB, because the RI immediate was 16 bits and no instruction loaded a
wide constant. V2 adds the U format:

\begin{lstlisting}[style=Cstyle]
  lui   r5, 0x123       # r5 = 0x123000
  addi  r5, r5, 0x456   # r5 = 0x123456
\end{lstlisting}

Any 32-bit address is therefore two instructions, and any load or store reaches
$\pm$32\,KiB around an arbitrary base. {\tt AUIPC} supplies the PC-relative form
for position-independent code.

\subsection{Conditional Branches}

The BR format carries two registers, so compare-and-branch is a single
instruction. V1 had only {\tt BZ} and {\tt BNZ}, which test one register against
zero; every other comparison compiled to a {\tt SUB} followed by a zero test.

Branches are PC-relative against the address of the instruction itself:

\[
\texttt{target} = \texttt{insn\_pc} + 4 + (\texttt{sign\_extend(offset)} \times 4)
\]

The offset is a 16-bit signed word offset, giving $\pm$128\,KiB of range. That
is a deliberate choice: with AUIPC available, a branch further than 128\,KiB is
two instructions, and spending bits on branch range would have cost either the
immediate width or the opcode width, both of which are load-bearing.

\subsection{Call, Return and Indirect Transfer}

{\tt CALL} writes the return address into the link register and jumps. {\tt RET}
jumps to the link register. {\tt JALR} is the indirect form, and is the
instruction that makes function pointers, virtual calls and callbacks
expressible --- V1 had no encoding for them at all, so plain C could not be
compiled against the ISA.

{\tt JALR} computes the target before writing the link, so a call through
{\tt lr} behaves correctly:

\begin{lstlisting}[style=Cstyle]
  jalr  ra, ra, 0       # t = pc+4; pc = ra + 0; ra = t  -> correct tail position
\end{lstlisting}

\subsection{Atomics and Memory Ordering}

{\tt LR.W} and {\tt SC.W} provide a load-reserved / store-conditional pair;
{\tt AMOSWAP.W} and {\tt AMOADD.W} provide unconditional atomic memory
operations. {\tt FENCE} is an ordering barrier with no operands.

The V2 atomic set is deliberately minimal --- four instructions, which is what a
working single-threaded and coarse-grained lock implementation needs. Bitwise
atomic operations (AND, OR, XOR) are not in V2. They are not forgotten: the
seven spare register-immediate slots and the four extension slots are where they
go when codegen demand is measured rather than guessed.

\subsection{Vector escape}
\label{sec:vector-escape}

Slots 0x3C--0x3F are reserved for extensions. The vector unit occupies this
space with a class field in the bits below the opcode, so vector instructions
are decoded from the same 64-entry table rather than sharing an opcode with a
scalar instruction.

In V1 the entire vector extension shared opcode {\tt 0xF} with {\tt RET}, and
was distinguished only by the value of the remaining bits, with
\texttt{0xF0000000} pinned as the sole {\tt RET} encoding. That is coherent as
written and fragile in practice: the whole vector extension depends on {\tt RET}
having exactly one encoding, and any future 32-bit pattern that collides becomes
a vector instruction silently. V1 also shared opcode {\tt 0xF} with the vector
opcode, so the scalar and vector decoders disagreed about what the same four
bits meant. V2 removes both by giving the vector unit its own opcode range.

\newpage

\newpage

\section{Exception and Interrupt Model}

BradISA V2 supports a vectored exception model with a fixed vector table at
the base of the address space.

\begin{table}[htbp]
\centering
\caption{Exception vector addresses}
\begin{tabular}{cll}
\toprule
Vector address & Exception       & Cause code \\
\midrule
0x00000000 & Reset / cold boot        & ---  \\
0x00000004 & Undefined instruction    & 1    \\
0x00000008 & Page fault               & 2    \\
0x0000000C & Unaligned access         & 3    \\
0x00000010 & Privilege violation      & 4    \\
0x00000014 & System call (ECALL)      & 5    \\
0x00000018 & Timer interrupt          & 6    \\
0x00000020 & External IRQ 0           & 7    \\
0x00000024 & External IRQ 1           & 7    \\
0x00000028 & External IRQ 2           & 7    \\
0x0000002C & External IRQ 3           & 7    \\
\bottomrule
\end{tabular}
\end{table}

\subsection{Machine Status Registers}
\begin{table}[htbp]
\centering
\caption{MSR index map}
\begin{tabular}{cll}
\toprule
Index & Mnemonic  & Description \\
\midrule
0     & STATUS    & Processor status flags \\
1     & CAUSE     & Cause of last exception \\
2     & EPC       & Exception program counter \\
3     & EAR       & Exception address register \\
4     & PAGE\_BASE & Page table base address \\
5     & TICK      & Cycle count timer \\
6     & CORE\_ID  & Core identifier \\
7     & CLUSTER\_ID & Cluster identifier \\
\bottomrule
\end{tabular}
\end{table}

\subsection{Status Register Flags}
\begin{table}[htbp]
\centering
\begin{tabular}{cll}
\toprule
Bit & Flag   & Description \\
\midrule
0   & IE     & Interrupt enable \\
1   & SVC    & Supervisor mode \\
2   & HALTED & Core halted \\
3   & WFE    & Wait for event \\
8   & EXC\_PEND & Exception pending \\
\bottomrule
\end{tabular}
\end{table}

\newpage

\section{Performance Characteristics}

\subsection{IPC Estimates}
\begin{itemize}[nosep]
  \item \textbf{Phoenix (10-stage OoO):} 2.2 IPC typical, 5.0 IPC peak
  \item \textbf{Falcon (5-stage in-order):} 0.9 IPC typical, 1.0 IPC peak
\end{itemize}

\subsection{Dhrystone Estimates}
\begin{itemize}[nosep]
  \item \textbf{Phoenix:} $\sim$5.5 DMIPS/MHz
  \item At 5.2\,GHz: $\sim$28,600 DMIPS
  \item \textbf{Falcon:} $\sim$2.2 DMIPS/MHz
  \item At 3.6\,GHz: $\sim$7,920 DMIPS
\end{itemize}

\subsection{Pipeline Stage Descriptions}

\subsubsection*{Phoenix (10-stage)}
\begin{enumerate}[nosep]
  \item \textbf{Fetch1} --- Instruction TLB access, next-PC selection
  \item \textbf{Fetch2} --- Instruction cache read
  \item \textbf{Decode} --- Instruction decode, register rename
  \item \textbf{Rename} --- Register map read, free-list access
  \item \textbf{Dispatch} --- Insert into issue queue / ROB
  \item \textbf{Issue} --- Wakeup, select, read physical register file
  \item \textbf{Execute} --- ALU / AGU / branch resolution
  \item \textbf{Memory} --- Data cache access (load/store)
  \item \textbf{Writeback} --- Result write to physical register file
  \item \textbf{Commit} --- Retire from ROB, update architectural state
\end{enumerate}

\subsubsection*{Falcon (5-stage)}
\begin{enumerate}[nosep]
  \item \textbf{Fetch1} --- Instruction fetch, PC increment
  \item \textbf{Fetch2} --- Instruction cache / memory read
  \item \textbf{Decode} --- Instruction decode, register file read
  \item \textbf{Execute} --- ALU / AGU / branch resolution, writeback
  \item \textbf{Writeback} --- Register file write
\end{enumerate}

\subsection{Vector Extension (VSET)}

BradISA V2 has no floating-point or SIMD facility; both are provided by the
BradISA V2 Vector Extension (VSET), specified in the companion document
\textbf{BradISA V2 Extensions}. VSET adds 32 $\times$ 256-bit vector
registers ({\tt v0}--{\tt v31}), elementwise integer/FP operations, vector
memory, reductions, and converts, encoded under the V2 extension-escape
slots {\tt 0x3C}--{\tt 0x3F} (\S\ref{sec:vector-escape}) so that no vector
instruction can share an encoding with {\tt RET}.
Enable via MSR~23 (VSTATUS.VE).

\newpage

\section*{Revision History}
\begin{tabular}{lll}
\toprule
Revision & Date & Changes \\
\midrule
1.0 & July 2026 & Initial specification --- Kinetic Gen1 \\
1.1 & July 2026 & Phoenix upgraded to octa-issue OoO (ROB 128), IPC 2.2/5.0, \\
    &           & DMIPS/MHz 5.5; Falcon DMIPS/MHz aligned to 2.2 \\
1.2 & July 2026 & Cross-reference BradISA V2 Vector Extension (VSET), \\
    &           & opcode {\tt 0xF}, VSTATUS MSR 23 \\
\bottomrule
\end{tabular}

\end{document}
